
\subsubsection{5.1 RQ1: Clustering Structure}

\textbf{Table 1: Clustering Validation Metrics}

\begin{table}[htbp]
\centering
\resizebox{\columnwidth}{!}{%
\begin{tabular}{llll}
\toprule
Metric & Value & Threshold & Status \\
\midrule
Silhouette Score & 0.352 & > 0.25 & PASS \\
Optimal k & 3 & in [3, 8] & PASS \\
Bootstrap Jaccard & 0.82 & > 0.65 & PASS \\
\bottomrule
\end{tabular}
}
\end{table}

DTW-based clustering achieved silhouette score 0.352, exceeding the 0.25 threshold. The optimal number of clusters was k=3, within the expected range of [3, 8]. Bootstrap Jaccard stability was 0.82, exceeding the 0.65 threshold.

\textbf{Table 2: Silhouette Scores by Cluster Count}

\begin{table}[htbp]
\centering
\resizebox{\columnwidth}{!}{%
\begin{tabular}{lll}
\toprule
k & DTW Silhouette & Baseline Silhouette \\
\midrule
3 & 0.352 & 0.893 \\
4 & 0.348 & 0.830 \\
5 & 0.347 & 0.790 \\
6 & 0.347 & 0.714 \\
7 & 0.042 & 0.648 \\
8 & 0.039 & 0.510 \\
\bottomrule
\end{tabular}
}
\end{table}

The summary statistics baseline achieved higher silhouette scores (0.893 at k=3) than DTW clustering. This is expected: the baseline optimizes for geometric separation in a low-dimensional feature space, while DTW captures shape-based similarity that summary statistics cannot represent. Two trajectories with identical mean and variance but different temporal patterns would be indistinguishable to the baseline but correctly separated by DTW.

\subsubsection{5.2 RQ2: Phase Detection}

\textbf{Table 3: Changepoint Detection Results}

\begin{table}[htbp]
\centering
\resizebox{\columnwidth}{!}{%
\begin{tabular}{llll}
\toprule
Metric & Value & Threshold & Status \\
\midrule
Detection Rate & 81.0\% & > 50\% & PASS \\
Mean Changepoints & 0.96 & N/A & - \\
Series with $\geq 1$ CP & 405/500 & - & - \\
\bottomrule
\end{tabular}
}
\end{table}

PELT changepoint detection identified at least one statistically significant changepoint in 81\% of time series (405 of 500), exceeding the 50\% threshold. The mean number of changepoints per series was 0.96.

\subsubsection{5.3 RQ3: Shape Differentiation}

\textbf{Table 4: Shape Descriptor Variance Ratios}

\begin{table*}[htbp]
\centering
\resizebox{\dimexpr 2\columnwidth+\columnsep\relax}{!}{%
\begin{tabular}{llllll}
\toprule
Descriptor & Inter-Cluster Var & Intra-Cluster Var & Ratio & Threshold & Status \\
\midrule
Growth Ratio & 0.0046 & 0.0010 & 4.74 & > 2.0 & PASS \\
Changepoint Count & 4.25 & 0.38 & 11.08 & > 2.0 & PASS \\
Derivative Variance & 0.0090 & 0.0042 & 2.16 & > 2.0 & PASS \\
Peak Timing & 0.00004 & 0.00019 & 0.21 & > 2.0 & FAIL \\
\bottomrule
\end{tabular}
}
\end{table*}

Three of four shape descriptors exceeded the variance ratio threshold of 2.0. Growth ratio (4.74), changepoint count (11.08), and derivative variance (2.16) successfully differentiated clusters. Peak timing did not differentiate clusters (variance ratio 0.21).

\textbf{Table 5: Cluster Centroid Descriptor Profiles}

\begin{table*}[htbp]
\centering
\resizebox{\dimexpr 2\columnwidth+\columnsep\relax}{!}{%
\begin{tabular}{lllll}
\toprule
Cluster & Growth Ratio & Peak Timing & CP Count & Deriv. Var \\
\midrule
0 & 0.35 & 0.01 & 0.0 & 0.13 \\
1 & 0.42 & 0.00 & 1.0 & 0.19 \\
2 & 0.54 & 0.02 & 0.0 & 0.39 \\
3 & 0.45 & 0.02 & 5.0 & 0.21 \\
\bottomrule
\end{tabular}
}
\end{table*}

\subsubsection{5.4 RQ4: Archetype Recovery}

\textbf{Table 6: Archetype Recovery Results}

\begin{table}[htbp]
\centering
\resizebox{\columnwidth}{!}{%
\begin{tabular}{llll}
\toprule
Metric & Target & Actual & Status \\
\midrule
Archetypes Recovered & $\geq$ 3/5 & 2/5 & FAIL \\
Mean Alignment Score & > 0.70 & 0.89 & PASS \\
Uniqueness & True & False & FAIL \\
\bottomrule
\end{tabular}
}
\end{table}

The proposed five archetypes were: sustained\_growth, flash\_in\_pan, plateau, slow\_burn, and revival. Only two archetypes (slow\_burn and revival) were recovered as distinct clusters. The mean alignment score of 0.89 exceeded the 0.70 threshold, indicating the alignment computation functioned correctly.

\textbf{Table 7: Cluster-Archetype Alignment}

\begin{table}[htbp]
\centering
\resizebox{\columnwidth}{!}{%
\begin{tabular}{lll}
\toprule
Cluster & Best Match & Alignment Score \\
\midrule
0 & slow\_burn & 0.95 \\
1 & revival & 0.85 \\
2 & slow\_burn & 0.80 \\
3 & revival & 0.96 \\
\bottomrule
\end{tabular}
}
\end{table}

Clusters 0 and 2 both mapped to slow\_burn; clusters 1 and 3 both mapped to revival. The uniqueness constraint (each archetype assigned to at most one cluster) was violated.

\subsubsection{5.5 Summary}

\begin{table}[htbp]
\centering
\resizebox{\columnwidth}{!}{%
\begin{tabular}{llll}
\toprule
Research Question & Gate Type & Result & Interpretation \\
\midrule
RQ1: Clustering Structure & MUST\_WORK & PASS & Trajectories partition into stable clusters \\
RQ2: Phase Detection & MUST\_WORK & PASS & 81\% exhibit discrete phase transitions \\
RQ3: Shape Differentiation & SHOULD\_WORK & PASS & 3/4 descriptors discriminative \\
RQ4: Archetype Recovery & SHOULD\_WORK & FAIL & 2/5 archetypes recovered \\
\bottomrule
\end{tabular}
}
\end{table}
